\documentclass[11pt]{article}

\usepackage[margin=1in]{geometry}
\usepackage{amsmath}
\usepackage{booktabs}
\usepackage{graphicx}
\usepackage{microtype}
\usepackage[hidelinks]{hyperref}

\title{Construction and Reconciliation of the Falcon 9 Infrasound Event Catalogue}
\author{}
\date{}

\begin{document}
\maketitle

\section{Purpose and provenance}

The Falcon 9 accident sequence was originally represented by two legacy
catalogue products containing 153 and 157 infrasound-event windows. Both were
derived from manual arrival picks made on the BCHH infrasound array using
Antelope \texttt{dbpick}, but the historical products used different event
constructions. We therefore treated both catalogues as provenance products and
reconstructed the event sequence directly from the preserved manual picks.
The reconstruction was designed to determine a reproducible event count
without selecting processing parameters merely because they reproduced either
legacy count.

The archived pick table contained 890 picks of all types. After limiting the
data to non-deleted infrasound phase-\(N\) picks within the accident interval,
384 picks remained. Of these, 66 were recorded on DD1, 162 on DD2, and 156 on
DD3. The archived channel names HD1--HD3 were mapped to the corresponding
current channel names DD1--DD3 without modifying the original pick table.
DD1 contained substantially fewer accepted picks than DD2 and DD3 and was
therefore treated as a supporting channel rather than as a requirement for
event acceptance.

\begin{table}[htbp]
\centering
\caption{Accepted manual infrasound picks used for catalogue reconstruction.}
\label{tab:accepted_picks}
\begin{tabular}{lrr}
\toprule
Channel & Accepted picks & Role in final assessment \\
\midrule
DD1 & 66  & Supporting channel; not allowed to veto DD2--DD3 detections \\
DD2 & 162 & Primary reference channel \\
DD3 & 156 & Primary corroborating channel \\
\bottomrule
\end{tabular}
\end{table}

\section{Reduced-time event association}

Raw inter-pick spacing was examined to document the legacy event construction,
but it was not used to select the production association window. The raw pick
times contain the acoustic travel time across the array; in particular, the
predicted DD1--DD2 differential travel time is approximately 0.090~s. A raw
0.1~s grouping window would therefore compensate partly for array geometry
rather than represent a physical separation between successive source events.

Each accepted pick was instead reduced to the DD2 reference time using

\begin{equation}
t_{i,\mathrm{red}}
=
t_i-
\frac{r_i-r_{\mathrm{DD2}}}{c_{\mathrm{eff}}},
\label{eq:reduced_time}
\end{equation}

where \(t_i\) is the observed pick time, \(r_i\) is the surveyed horizontal
distance from SLC-40 to sensor \(i\), and
\(c_{\mathrm{eff}}=350.977~\mathrm{m\,s^{-1}}\) is the meteorologically
estimated effective sound speed along the source-to-array path. The resulting
corrections were 0.090107~s for DD1, 0 for DD2, and 0.012867~s for DD3. These
corrections were used only for pick association; the continuous waveform time
axes were not altered.

Reduced picks were sorted chronologically. Beginning with the earliest
unassigned pick, all remaining picks within a following association window
were considered a candidate group. At most one pick per channel was retained,
and a candidate event required picks from at least two distinct infrasound
channels. Picks that did not satisfy this requirement remained unassociated
rather than being forced into an event.

Association windows from 0.008 to 0.060~s were tested. The reconstructed event
count increased from 128 at 0.008~s to 154 at 0.040~s and remained unchanged
through 0.060~s (Table~\ref{tab:window_sensitivity}). We selected 0.040~s as
the smallest tested value on this stable plateau. At this value, 371 of the
384 accepted picks were associated into 154 candidates: 63 candidates had
three-channel support and 91 had two-channel support. The median reduced pick
span was 0.004867~s, and its 95th percentile was 0.017226~s.

\begin{table}[htbp]
\centering
\caption{Sensitivity of the reduced-time catalogue to the association window.}
\label{tab:window_sensitivity}
\begin{tabular}{rrrrrr}
\toprule
Window & Events & Associated & Unassociated & Three-channel & Two-channel \\
(s) & & picks & picks & events & events \\
\midrule
0.008 & 128 & 303 & 81 & 47 & 81 \\
0.012 & 144 & 347 & 37 & 59 & 85 \\
0.016 & 149 & 357 & 27 & 59 & 90 \\
0.020 & 151 & 363 & 21 & 61 & 90 \\
0.024 & 152 & 366 & 18 & 62 & 90 \\
0.030 & 153 & 368 & 16 & 62 & 91 \\
\textbf{0.040} & \textbf{154} & \textbf{371} & \textbf{13} & \textbf{63} & \textbf{91} \\
0.050 & 154 & 371 & 13 & 63 & 91 \\
0.060 & 154 & 371 & 13 & 63 & 91 \\
\bottomrule
\end{tabular}
\end{table}

The 0.040~s window should be interpreted as an operational association
tolerance rather than as a direct sound-speed constraint. Across the longest
source-distance difference, plausible apparent speeds of approximately
280--440~\(\mathrm{m\,s^{-1}}\) would produce residual delays no larger than
about 0.023~s after reduction. The additional tolerance accommodates manual
pick uncertainty, differences in waveform onset among sensors, small source
position variations, and propagation-model error.

\section{Targeted candidate review and final event count}

Only one candidate was added when the association window increased from 0.030
to 0.040~s. Candidate 137, at approximately 13:22:29.928 UTC, was supported by
DD2 and DD3 with a reduced pick span of 0.035866~s. After waveform alignment,
its best DD2--DD3 correlation was approximately 0.879 and its robust stack
peak signal-to-noise ratio (SNR) was approximately 8.0. The coherent waveform
supports retaining this threshold-sensitive event.

One of the 154 candidates, candidate 148 at approximately 13:24:50.831 UTC,
was supported only by DD1 and DD2. It had no DD3 pick or coherent DD3 waveform
counterpart; the best DD2--DD3 correlation was only approximately 0.092.
Because DD1 was the least reliable array element, DD1--DD2 support alone was
not considered sufficient to establish an infrasound event. Candidate 148 was
therefore classified as ambiguous and excluded. The production detection
catalogue consequently contains

\begin{equation}
154\ \text{reduced-time candidates}
-1\ \text{ambiguous candidate}
=153\ \text{detected events}.
\end{equation}

This result supports the historical count of 153, but it is not obtained by
tuning the association window to that number. Instead, it follows from a
stable 154-candidate association result and an explicit waveform-based
decision concerning the sole DD1--DD2-only candidate.

Direct assignment of the accepted picks to the legacy 157-event windows also
showed why the larger historical catalogue should not be interpreted as 157
independently confirmed events. Only 146 of the 157 windows had support from
at least two channels, 11 had fewer than two supporting channels, and three
contained no accepted infrasound picks. Four adjacent window pairs touched or
overlapped. Comparison of the two legacy products identified two possible
153-to-157 splits and no convincing merges.

\section{Waveform-quality subsets}

Event detection was separated from suitability for quantitative waveform and
array analysis. For every retained event, subsequent processing measured
robust peak SNR, the signal-window to noise-window standard-deviation ratio,
shifted DD2--DD3 correlation and residual lag, pressure amplitude, pulse
morphology, and three-channel array coherence where DD1 was usable.

The robust noise scale was defined using the median absolute deviation,

\begin{equation}
\sigma_{\mathrm{robust}}
=
1.4826\,\mathrm{median}
\left(
\left|p-\mathrm{median}(p)\right|
\right),
\end{equation}

and robust peak SNR was calculated relative to this pre-event noise estimate.
Of the 153 detected events, 138 had robust peak SNR of at least 3 and shifted
DD2--DD3 correlation of at least 0.5. A core subset of 107 events had robust
peak SNR of at least 5 and DD2--DD3 correlation of at least 0.7. Fifteen
events remained valid DD2--DD3-supported detections but fell below the usable
waveform thresholds.

\begin{table}[htbp]
\centering
\caption{Operational catalogue and waveform-quality subsets. Thresholds are
descriptive quality divisions rather than natural event classes.}
\label{tab:quality_subsets}
\begin{tabular}{p{0.22\linewidth}p{0.58\linewidth}r}
\toprule
Subset & Criteria & Events \\
\midrule
Detection catalogue & DD2 and DD3 reduced-time pick support after targeted review & 153 \\
Usable waveform subset & Robust peak SNR $\geq 3$ and DD2--DD3 correlation $\geq 0.5$ & 138 \\
Core waveform subset & Robust peak SNR $\geq 5$ and DD2--DD3 correlation $\geq 0.7$ & 107 \\
Marginal detections & Retained detections below the usable waveform thresholds & 15 \\
\bottomrule
\end{tabular}
\end{table}

The absolute signal-window root-mean-square (RMS) or standard deviation was
not used as an additional catalogue-defining cutoff. Absolute RMS is primarily
an amplitude measure, while the signal-to-noise RMS ratio depends strongly on
the selected window duration. A short pressure impulse can have high peak SNR
but only moderate windowed RMS because the pulse occupies a small fraction of
the measurement interval. The RMS or standard-deviation ratio should therefore
be retained as a continuous diagnostic and sensitivity measure rather than as
a binary event criterion.

\section{Recommended presentation in the paper}

\subsection{Main text}

No new main-text figure is essential if the paper already contains a
continuous-waveform or event-chronology figure showing the 153-event sequence.
The principal catalogue-construction results can be reported concisely in the
Methods and Results:

\begin{itemize}
\item 384 accepted manual infrasound picks;
\item differential-time reduction to DD2 using 350.977~\(\mathrm{m\,s^{-1}}\);
\item a 0.040~s association window and a two-channel minimum;
\item 154 candidates followed by one explicit ambiguous-event exclusion;
\item a final 153-event detection catalogue;
\item 138 usable and 107 core waveform-quality events.
\end{itemize}

Table~\ref{tab:quality_subsets} is the only compact table that may be useful in
the main paper, although these counts could also be stated in prose. The full
association-window table is methodological detail and need not occupy main-text
space.

\subsection{Supplementary material}

Two supplementary figures would provide the strongest audit trail:

\begin{enumerate}
\item \textbf{Reduced-time association sensitivity.} Event count, associated
pick count, and preferably the median or 95th-percentile corrected pick span
versus association window. This demonstrates the stable 0.040--0.060~s
plateau. The existing event-count figure is adequate, but a second panel for
pick-span statistics would make it more informative.
\item \textbf{Targeted waveform decisions.} Baseline-corrected, reduced-time
DD2 and DD3 overlays for candidate 137 and candidate 148. This directly
documents why the threshold-sensitive candidate was retained and the
DD1--DD2-only candidate was excluded.
\end{enumerate}

The following should be supplied electronically:

\begin{itemize}
\item the complete 153-event catalogue;
\item the 154-candidate decision table, including the excluded candidate and
its rationale;
\item the reduced-pick assignment table;
\item the association-window sensitivity table;
\item per-event waveform-quality measurements and tier assignments.
\end{itemize}

The raw global-gap histogram, same-channel-gap histogram, raw-time clustering
sweep, 157-window pick raster, and complete 153-versus-157 matching tables are
valuable provenance products but are not essential publication figures. They
can remain in the reproducibility archive or an extended supplement. In
particular, the raw-time threshold plots should not be emphasized because they
combine source recurrence with inter-sensor propagation delays.

\end{document}
